% ============================================================
%  Chapter 2 — Meet the Brain: The Processor
%  Inside a Smartphone: The Tiny World in Your Hand
% ============================================================

\chapter{Meet the Brain: The Processor}

\chapteropening{%
  Everything your phone does — every tap, swipe, photo, and song —
  starts with one tiny chip doing billions of calculations every second.
}
\chapterbadge{Mission: Follow an Instruction Through the Brain!}

% --- Opening ---
\section*{Meet Chip}

Dot led Sam deeper into the phone, past glowing wires and humming circuits,
until they reached the heart of the device — a flat, shiny square no bigger than a fingernail.

A tiny character with sparkling eyes and a crisp silver jacket waved at them.

\character{Chip}{Welcome! I am Chip — the processor.
  Everything you do on this phone passes through me first.
  When you tap the screen, I get the message.
  When you take a photo, I handle it.
  Even reading this book on a screen would go through me.}

\sectionrule

% --- What a processor does ---
\section*{What Does a Processor Do?}

A \term{processor} (also called a CPU — Central Processing Unit) is the brain of the smartphone.
Its job is to follow \term{instructions}.

Every action you take on a phone is turned into a long list of instructions.
The processor reads each instruction and carries it out, one after another,
billions of times every second.

\begin{picturethis}
  The processor is like a chef in a very busy kitchen.
  Orders keep arriving: chop this, heat that, plate this next.
  The chef follows each instruction in order, very fast.
  The processor does the same thing — except it handles billions
  of instructions in the time it takes you to blink.
\end{picturethis}

\sectionrule

% --- Transistors ---
\section*{Billions of Tiny Switches}

Inside a processor are billions of microscopic switches called \term{transistors}.
Each transistor can be either \textbf{on} (1) or \textbf{off} (0).
These ones and zeros are called \term{binary code} — the language all computers use.
Groups of ones and zeros can represent numbers, words, colours, and sounds.

\begin{funfact}
  A modern smartphone processor contains more than \textbf{10 billion transistors}.
  Each one is incredibly tiny. If you lined transistors up across a single
  human hair, thousands would fit.
\end{funfact}

\sectionrule

% --- Speed ---
\section*{Why Speed Matters}

Processor speed is measured in \term{gigahertz} (GHz).
One gigahertz means one billion steps per second.

\begin{quickmap}
  \textbf{What happens when you tap a photo to open it:}
  \begin{enumerate}[label=\textbf{\arabic*.}]
    \item The touchscreen detects your tap and sends a signal.
    \item The processor receives the signal (an instruction).
    \item The processor fetches the photo data from storage.
    \item The processor unpacks the image data into pixels.
    \item The processor sends the image to the screen.
    \item All of this happens in a fraction of a second.
  \end{enumerate}
\end{quickmap}

\begin{diagram}[The processor connects to all other parts of the phone.]
  \begin{tikzpicture}[>=Stealth, node distance=1.8cm]
    \node[draw=TechPurple,fill=TechPurple!20,rounded corners,
          minimum width=2.4cm,minimum height=1.2cm,line width=1.2pt,
          align=center,font=\bfseries] (cpu) {Processor\\(Chip)};
    \node[draw=SoftPurple!70,fill=SoftPurple!10,rounded corners,
          minimum width=2cm,minimum height=0.8cm,above=of cpu] (ram) {RAM};
    \node[draw=LeafGreen!70,fill=LeafGreen!10,rounded corners,
          minimum width=2cm,minimum height=0.8cm,right=of cpu] (store) {Storage};
    \node[draw=WarmOrange!70,fill=WarmOrange!10,rounded corners,
          minimum width=2cm,minimum height=0.8cm,below=of cpu] (screen) {Screen};
    \node[draw=CoralRed!70,fill=CoralRed!10,rounded corners,
          minimum width=2cm,minimum height=0.8cm,left=of cpu] (touch) {Touchscreen};
    \draw[<->,thick,TechPurple] (cpu) -- (ram);
    \draw[<->,thick,TechPurple] (cpu) -- (store);
    \draw[<->,thick,TechPurple] (cpu) -- (screen);
    \draw[<->,thick,TechPurple] (cpu) -- (touch);
  \end{tikzpicture}
\end{diagram}

\sectionrule

\begin{newwords}
  \begin{description}[labelwidth=3.5cm, leftmargin=3.7cm]
    \item[\term{Processor}]    \emph{(say: PROSS-ess-er)} — The brain chip that carries out all instructions.
    \item[\term{Transistor}]   \emph{(say: tran-ZISS-ter)} — A microscopic on/off switch inside a processor.
    \item[\term{Binary code}]  \emph{(say: BY-nuh-ree)} — The language of ones and zeros that computers use.
    \item[\term{Gigahertz}]    \emph{(say: GIG-uh-hurts)} — A unit of speed: one billion steps per second.
    \item[\term{Instruction}]  \emph{(say: in-STRUK-shun)} — A single command the processor must carry out.
  \end{description}
\end{newwords}

\sectionrule

\begin{tryit}
  \textbf{Be the processor!}

  Ask a friend to give you instructions one at a time
  (e.g., ``Pick up the pencil'', ``Move it to the table'', ``Put it down'').
  Follow each instruction \emph{exactly} as given.
  What happens if an instruction is unclear?

  \emph{This is why programmers must write instructions very carefully.}
\end{tryit}

\sectionrule

\begin{bigidea}
  The \textbf{processor} is the brain of the phone.
  It follows billions of instructions every second using tiny transistors
  and controls everything else inside the phone.
\end{bigidea}

\character{Chip}{I need somewhere to hold all the information I am working on.
  That is where my friend Memo comes in. On to Chapter 3!}
